% Results shared with the Paper; written by src/build_report_fragments.py.
% Generated by src/build_report_fragments.py from outputs/paper_audit_20261024.
% Do not edit by hand; rerun the generator after rebuilding the evidence.
\newcommand{\CalBSIVRMSE}{131.58}
\newcommand{\CalBSIVRMSELegacy}{97.21}
\newcommand{\CalHIVRMSE}{68.15}
\newcommand{\CalHIVRMSELegacy}{49.49}
\newcommand{\CalBIVRMSE}{65.24}
\newcommand{\CalBIVRMSELegacy}{49.41}
\newcommand{\CalWIVRMSE}{52.83}
\newcommand{\CalWIVRMSELegacy}{49.33}
\newcommand{\HawkesBranching}{0.825}
\newcommand{\HawkesBeta}{10.27}
\newcommand{\HawkesHalfLifeDays}{24.7}
\newcommand{\HawkesMeanIntensity}{1.65}
\newcommand{\BatesLambda}{1.72}
\newcommand{\HestonKappa}{2.93}
\newcommand{\FellerRelaxedLossRatio}{68}
\newcommand{\KappaFlatLow}{2}
\newcommand{\KappaFlatHigh}{4}
\newcommand{\CarryPriceRMSEZero}{0.49}
\newcommand{\CarryPriceRMSEFee}{0.75}
\newcommand{\FuturesYieldLow}{0.26}
\newcommand{\FuturesYieldHigh}{0.49}
\newcommand{\AmericanPremiumMax}{2.3e-10}
\newcommand{\ValBSMedian}{56.32}
\newcommand{\ValBSPTen}{4.31}
\newcommand{\ValBSPNinety}{160.19}
\newcommand{\ValBSMean}{71.94}
\newcommand{\ValBSNegative}{7.4}
\newcommand{\ValBSExplicit}{10.75}
\newcommand{\ValBSExplicitRiskFree}{11.81}
\newcommand{\ValHMedian}{57.62}
\newcommand{\ValHPTen}{0.13}
\newcommand{\ValHPNinety}{159.15}
\newcommand{\ValHMean}{71.63}
\newcommand{\ValHNegative}{9.9}
\newcommand{\ValHExplicit}{10.90}
\newcommand{\ValHExplicitRiskFree}{11.98}
\newcommand{\ValBMedian}{59.48}
\newcommand{\ValBPTen}{2.09}
\newcommand{\ValBPNinety}{156.74}
\newcommand{\ValBMean}{72.05}
\newcommand{\ValBNegative}{8.9}
\newcommand{\ValBExplicit}{11.08}
\newcommand{\ValBExplicitRiskFree}{12.19}
\newcommand{\ValWMedian}{57.26}
\newcommand{\ValWPTen}{1.26}
\newcommand{\ValWPNinety}{155.88}
\newcommand{\ValWMean}{73.11}
\newcommand{\ValWNegative}{9.3}
\newcommand{\ValWExplicit}{10.90}
\newcommand{\ValWExplicitRiskFree}{11.97}
\newcommand{\BarBS}{37.87}
\newcommand{\BarH}{38.75}
\newcommand{\BarB}{40.00}
\newcommand{\BarW}{38.51}
\newcommand{\ValMedianLow}{56.32}
\newcommand{\ValMedianHigh}{59.48}
\newcommand{\ValExplicitLow}{10.75}
\newcommand{\ValExplicitHigh}{11.08}
\newcommand{\ValPNinetyLow}{155.88}
\newcommand{\ValPNinetyHigh}{160.19}
\newcommand{\TerminalShare}{83}
\newcommand{\NumericalMedianLow}{56.24}
\newcommand{\NumericalMedianHigh}{59.99}
\newcommand{\NumericalHawkesLow}{56.26}
\newcommand{\NumericalHawkesHigh}{58.42}
\newcommand{\AttLegacyBS}{60.10}
\newcommand{\AttLegacyH}{59.77}
\newcommand{\AttLegacyB}{59.12}
\newcommand{\AttLegacyW}{59.82}
\newcommand{\AttClockBS}{57.81}
\newcommand{\AttClockH}{58.65}
\newcommand{\AttClockB}{59.34}
\newcommand{\AttClockW}{58.49}
\newcommand{\AttAnchorBS}{57.11}
\newcommand{\AttAnchorH}{57.95}
\newcommand{\AttAnchorB}{58.59}
\newcommand{\AttAnchorW}{57.74}
\newcommand{\AttCarryBS}{56.32}
\newcommand{\AttCarryH}{57.62}
\newcommand{\AttCarryB}{59.48}
\newcommand{\AttCarryW}{57.26}
\newcommand{\CuBSGold}{12.060}
\newcommand{\CuBSCombined}{14.368}
\newcommand{\CuBSDelta}{2.308}
\newcommand{\CuBSSE}{17.9}
\newcommand{\CuHGold}{11.981}
\newcommand{\CuHCombined}{14.285}
\newcommand{\CuHDelta}{2.304}
\newcommand{\CuHSE}{28.5}
\newcommand{\CuBGold}{12.062}
\newcommand{\CuBCombined}{14.379}
\newcommand{\CuBDelta}{2.317}
\newcommand{\CuBSE}{31.5}
\newcommand{\CuWGold}{12.154}
\newcommand{\CuWCombined}{14.484}
\newcommand{\CuWDelta}{2.330}
\newcommand{\CuWSE}{28.1}
\newcommand{\CuDeltaLow}{2.30}
\newcommand{\CuDeltaHigh}{2.33}
\newcommand{\TerminalWACC}{8.445}
\newcommand{\TerminalMultiple}{13.14}
\newcommand{\WACCClippedShare}{0.00}
\newcommand{\HorizonYears}{4.33}
\newcommand{\StubDays}{28}
\newcommand{\CovBSNinetyFive}{27/29}
\newcommand{\CovBSNinetyFiveP}{0.655}
\newcommand{\CovBSEighty}{21/29}
\newcommand{\CovBSEightyP}{0.350}
\newcommand{\CovHNinetyFive}{27/29}
\newcommand{\CovHNinetyFiveP}{0.655}
\newcommand{\CovHEighty}{20/29}
\newcommand{\CovHEightyP}{0.160}
\newcommand{\CovBNinetyFive}{25/29}
\newcommand{\CovBNinetyFiveP}{0.055}
\newcommand{\CovBEighty}{17/29}
\newcommand{\CovBEightyP}{0.008}
\newcommand{\CovWNinetyFive}{25/29}
\newcommand{\CovWNinetyFiveP}{0.055}
\newcommand{\CovWEighty}{16/29}
\newcommand{\CovWEightyP}{0.004}
\newcommand{\OosBSRMSE}{163.35}
\newcommand{\OosBSRMean}{-2.33}
\newcommand{\OosBSRPersistence}{-423.92}
\newcommand{\OosBSBeatsPersistence}{0}
\newcommand{\OosBSDMPersistenceP}{0.000}
\newcommand{\OosHRMSE}{89.62}
\newcommand{\OosHRMean}{69.17}
\newcommand{\OosHRPersistence}{-59.00}
\newcommand{\OosHBeatsPersistence}{4}
\newcommand{\OosHDMPersistenceP}{0.001}
\newcommand{\OosBRMSE}{68.39}
\newcommand{\OosBRMean}{80.62}
\newcommand{\OosBRPersistence}{-17.47}
\newcommand{\OosBBeatsPersistence}{14}
\newcommand{\OosBDMPersistenceP}{0.088}
\newcommand{\OosWRMSE}{64.16}
\newcommand{\OosWRMean}{83.03}
\newcommand{\OosWRPersistence}{-1.84}
\newcommand{\OosWBeatsPersistence}{13}
\newcommand{\OosWDMPersistenceP}{0.772}
\newcommand{\OosDMHawkesVsBSP}{0.000}
\newcommand{\OosDMHawkesVsHP}{0.000}
\newcommand{\OosDMHawkesVsBP}{0.165}
\newcommand{\OosDates}{30}
\newcommand{\OosObservations}{4,667}
\newcommand{\OosPersistenceObservations}{4,026}
\newcommand{\OosPersistenceRMSE}{58.31}
\newcommand{\FigOosRmse}{../img/publication/7d3e5597f15ab728_oos_rmse.png}
\newcommand{\FigOosRsquared}{../img/publication/19d304216f437e9e_oos_r2.png}
\newcommand{\FigFigValPrimary}{../img/publication/0e6f7e2f5b6abc45_fig_val_primary.png}
\newcommand{\FigFigValMulti}{../img/publication/893e0badbf0a67cf_fig_val_multi.png}
\newcommand{\FigFigTerminalShare}{../img/publication/36af633f72f0ca22_fig_terminal_share.png}
\newcommand{\FigFigCorrectionAttribution}{../img/publication/6d67719615b71840_fig_correction_attribution.png}
\newcommand{\FigBlackScholesResidualHeatmap}{../img/publication/612c9067fff153cc_black_scholes_residual_heatmap.png}
\newcommand{\FigHestonResidualHeatmap}{../img/publication/ba2dd686e32e5e10_heston_residual_heatmap.png}
\newcommand{\FigBatesResidualHeatmap}{../img/publication/917f624afe7e339f_bates_residual_heatmap.png}
\newcommand{\FigBatesHawkesResidualHeatmap}{../img/publication/38343b91b0c0ad53_bates_hawkes_residual_heatmap.png}
\newcommand{\FigArfimaFdSilver}{../img/publication/a47ac9bb3870cd03_arfima_fd_silver.png}
\newcommand{\GphGoldAbs}{0.452}
\newcommand{\GphSilverAbs}{0.743}


\chapter*{Discoveries at a Glance}
\addcontentsline{toc}{chapter}{Discoveries at a Glance}
\markboth{Discoveries at a Glance}{Discoveries at a Glance}

\begin{tcolorbox}[
  title={Discovery 1: the gold-price model matters less than numerical noise},
  colback=sflightblue,
  colframe=sfblue,
  boxrule=0.7pt,
  arc=1mm,
  left=3mm,
  right=3mm
]
\small
Holding production, costs and discount-rate shocks fixed, the four gold-price
laws give median operating proxies between \textbf{USD~\ValMedianLow{} and
\ValMedianHigh~billion}. Changing only the number of paths or the seed moves a
single model's median across \textbf{USD~\NumericalMedianLow--\NumericalMedianHigh~billion}.
The spread between models is therefore no larger than the Monte Carlo
uncertainty of one model, and the medians do not rank the models.
\end{tcolorbox}

\begin{tcolorbox}[
  title={Discovery 2: the terminal value carries the result},
  colback=sflightblue,
  colframe=sfblue,
  boxrule=0.7pt,
  arc=1mm,
  left=3mm,
  right=3mm
]
\small
Flows to 31 December 2030 are worth \textbf{USD~\ValExplicitLow--\ValExplicitHigh~billion};
the perpetuity after 2030 carries about \textbf{\TerminalShare\%} of the mean proxy.
The terminal assumptions, not the gold model, decide the level.
\end{tcolorbox}

\begin{tcolorbox}[
  title={Discovery 3: no option model beats the persistence of the IV surface},
  colback=sflightblue,
  colframe=sfblue,
  boxrule=0.7pt,
  arc=1mm,
  left=3mm,
  right=3mm
]
\small
In 30 rolling repricing tests, Bates--Hawkes has the lowest error among the
four models (\OosWRMSE~bp), but carrying the previous day's implied-volatility
surface forward is still more accurate: the out-of-sample $R^2$ against
persistence is \OosWRPersistence\% for Bates--Hawkes and lower for the others.
\end{tcolorbox}

\noindent\textbf{What these numbers are.} They are operating-margin proxies,
not enterprise or equity values: the equity bridge (debt, cash, minorities,
other claims and the diluted share count) is not reconciled, so no value per
share is reported and no comparison with Barrick's share price is made. The
results are research sensitivities, not fair values, target prices or
investment advice.

\noindent\textbf{What the four models do---and do not do.} Black--Scholes/GBM,
Heston, Bates--Poisson and Full Bates--Hawkes govern only the \emph{gold-price
layer}. They are calibrated to 64 GLD calls of 2 September 2026 with the
trust's 0.40\% sponsor fee as carry. They do not generate mine production or
costs: those come from Team~4's operating forecasts, identical in every run.
The unified DCF and the stochastic WACC shocks are common as well. Only flows
after the 2 September 2026 valuation time are valued; Barrick's Q1--Q2 2026
actuals have already been realised and are excluded.

\begin{table}[H]
\centering
\caption{Median operating proxy under each gold-price law, USD billion. Explicit: flows to 2030 without terminal value.}
\label{tab:discoveries-medians}
\small
\begin{tabular}{@{}lrrr@{}}
\toprule
Gold engine & Median & P10--P90 & Explicit \\
\midrule
Black--Scholes/GBM & \ValBSMedian & \ValBSPTen--\ValBSPNinety & \ValBSExplicit \\
Heston & \ValHMedian & \ValHPTen--\ValHPNinety & \ValHExplicit \\
Bates--Poisson & \ValBMedian & \ValBPTen--\ValBPNinety & \ValBExplicit \\
Full Bates--Hawkes & \ValWMedian & \ValWPTen--\ValWPNinety & \ValWExplicit \\
\bottomrule
\end{tabular}
\end{table}

\begin{figure}[H]
  \centering
  \includegraphics[width=0.88\textwidth]{\FigFigValMulti}
  \caption{Operating-proxy distributions under the four gold-price laws on the
  dated clock, 8,192 paths each, with common operations, WACC shocks and
  signed terminal closure. No share count, equity bridge or market reference
  is used.}
  \label{fig:discoveries-multimodel}
\end{figure}

\noindent\textbf{How the evidence was assembled.} The comparison is the
endpoint of an integrated chain. Team~8 supplies option-implied gold-price
dynamics; Team~4 supplies the production and unit-cost layer; Team~5 maps the
resulting component margin through cash-flow and discounting identities;
Team~3 supplies simulation diagnostics. The remaining teams provide
probability, econometric, historical-volatility and stochastic-process
foundations. Integration required one owner for each variable, common units
and dates, and separate code, factual, editorial and visual reviews.

The chain separates calibration from forecasting. Option parameters are
estimated under the risk-neutral measure; the scenarios start from the gold
spot of USD~4,391.32/oz at the valuation time and drift at the Treasury
forward rate, which as a physical expectation assumes no gold risk premium.
The option-implied volatility shape is carried over as a stated assumption,
not a validated physical law. Section~\ref{sec:physical-measure-bridge} reports the separate physical
forecasting evidence.

\begin{table}[H]
\centering
\caption{Evidence behind the valuation layer.}
\label{tab:discoveries-evidence}
\small
\renewcommand{\arraystretch}{1.13}
\begin{tabularx}{\textwidth}{@{}p{2.65cm}X X@{}}
\toprule
\textbf{Evidence block} & \textbf{Dated observation} &
\textbf{Permitted interpretation} \\
\midrule
Barrick equity, \texttt{B} & 86 valid daily unadjusted-close bars from
27 April to 27 August 2026 (LSE run \texttt{20260828T070000Z-lse-nss});
cumulative change 16.20\%, annualised log-return volatility 51.23\%, maximum
drawdown $-26.10\%$. & Short-window market diagnostic only; not a long-run
return or volatility estimate. \\
GLD reference series & 1,425 daily bars from 4 January 2021 to 27 August 2026;
total return 131.06\%, annualised volatility 18.21\%, maximum drawdown
$-26.23\%$. & ETF price in USD/share; neither physical gold in USD/oz nor
Barrick's realised selling price. \\
GLD surface & 605 eligible calls; CC64 selects 64 actual contracts across eight
maturities and 20 strikes on 2 September 2026; same-date 14-tenor Treasury
curve; carry $q=0.40\%$. & Option-implied calibration input; no physical
drift is identified. \\
Path engines & Daily full-truncation schemes with exact Hawkes intensity
integration; Monte Carlo call prices within about two batch standard errors of
COS prices for all four models. & Numerical correctness of the simulator, not
evidence that the law is physically right. \\
Integrated experiment & 8,192 paths; four gold engines; identical operating
contract, DCF and WACC shocks. & A controlled attribution of distributional
differences to the gold-price engine. \\
\bottomrule
\end{tabularx}
\end{table}

\noindent\textbf{Who produced the study.} Research is the Starting Finance
Club PoliTo area in which members collaborate on quantitative, economic,
financial and methodological projects. The Barrick programme joined eight
autonomous teams with backgrounds spanning management engineering,
mathematics, data science, computation and quantitative finance. Common
foundations and coordination rules allowed specialised work to remain deep
while still being reviewed and reused across teams. The unified study keeps
their authorship and source evidence, and treats integration as a distinct
phase with its own code, writing, technical, editorial and visual reviews.

The chapters that follow reconstruct the company and data perimeter, the
mathematical and empirical model layers, the common operating and DCF
contracts, the four-model experiment and the limits that keep a conditional
signal from becoming an investment conclusion.
